% =====================================================================
%  ANNEXE F — JOURNAL DE PROGRESSION (à imprimer)
% =====================================================================
\chapter{Votre journal de progression}

Cette annexe est faite pour être \textbf{imprimée, remplie à la main, et conservée}.
Un journal de progression remplit trois fonctions : il rend l'effort visible, il
transforme une intention en rendez-vous, et il fournit --- dans les moments de doute ---
la preuve que vous avancez. Au mois 8, quand la motivation fléchira, vous relirez les
pages des mois 2 et 3.

\begin{regle}[La seule règle du journal]
\textbf{On remplit le journal le jour même, jamais le lendemain.} Deux minutes suffisent :
trois cases à cocher, un chiffre à écrire, une phrase à noter. Un journal tenu une fois par
semaine devient un journal abandonné en un mois.
\end{regle}

\section{Le contrat personnel}
À remplir le jour 1, à signer, et à relire chaque premier dimanche du mois.

\begin{tableau}[Mon contrat d'apprentissage]
\begin{tabularx}{\linewidth}{@{}p{7.4cm}X@{}}
\entete{Question} & \entete{Ma réponse} \tabularnewline
\hline
Pourquoi est-ce que je veux parler anglais ? (la vraie raison, pas la raison polie) & \rule{0pt}{2.2cm}\trou[6.6cm] \tabularnewline
Qu'est-ce que cela changera concrètement dans ma vie ? & \rule{0pt}{2.0cm}\trou[6.6cm] \tabularnewline
Combien de minutes par jour est-ce que je m'engage à y consacrer ? & \rule{0pt}{1.2cm}\trou[6.6cm] \tabularnewline
À quelle heure exactement ? (créez le rendez-vous) & \rule{0pt}{1.2cm}\trou[6.6cm] \tabularnewline
Où ? (le lieu compte autant que l'heure) & \rule{0pt}{1.2cm}\trou[6.6cm] \tabularnewline
Quelle est la plus petite version possible de ma séance, les jours difficiles ? & \rule{0pt}{1.6cm}\trou[6.6cm] \tabularnewline
Qui sera au courant de mon engagement ? & \rule{0pt}{1.2cm}\trou[6.6cm] \tabularnewline
Quelle récompense si je tiens 90 jours ? & \rule{0pt}{1.2cm}\trou[6.6cm] \tabularnewline
Ma date de début : & \rule{0pt}{1.0cm}\trou[4.0cm] \tabularnewline
Ma date de fin (365 jours plus tard) : & \rule{0pt}{1.0cm}\trou[4.0cm] \tabularnewline
Signature : & \rule{0pt}{1.2cm}\trou[6.6cm] \tabularnewline
\end{tabularx}
\end{tableau}

\begin{astuce}[Le seuil minimal des mauvais jours]
Écrivez dans votre contrat une « plus petite version possible » ridiculement facile :
\emph{5 minutes --- 5 cartes Anki et une phrase à voix haute}. Le but n'est pas de
progresser ces jours-là, mais de \textbf{ne pas casser la chaîne}. Une chaîne de 300 jours
avec vingt séances de 5 minutes vaut infiniment mieux qu'une chaîne brisée trois fois.
\end{astuce}

\section{La grille hebdomadaire}
À photocopier 52 fois. Une ligne par jour, une case par activité. Une case remplie vaut
mieux qu'une case parfaite.

\newcommand{\case}{\fbox{\hspace{1.0em}\rule[-0.35em]{0pt}{1.05em}}}
\begin{tableau}[Semaine du \rule{2cm}{0.4pt} au \rule{2cm}{0.4pt}]
\footnotesize
\begin{tabularx}{\linewidth}{@{}p{1.5cm}p{1.0cm}p{1.0cm}p{1.0cm}p{1.0cm}p{1.0cm}p{1.0cm}p{1.0cm}X@{}}
\entete{Jour} & \entete{15 min Anki} & \entete{15 min grammaire} & \entete{20 min écoute} & \entete{15 min lecture} & \entete{10 min écrit} & \entete{20 min oral} & \case & \entete{Note du jour} \tabularnewline
\hline
Lundi & \case & \case & \case & \case & \case & \case & \case & \tabularnewline
Mardi & \case & \case & \case & \case & \case & \case & \case & \tabularnewline
Mercredi & \case & \case & \case & \case & \case & \case & \case & \tabularnewline
Jeudi & \case & \case & \case & \case & \case & \case & \case & \tabularnewline
Vendredi & \case & \case & \case & \case & \case & \case & \case & \tabularnewline
Samedi & \case & \case & \case & \case & \case & \case & \case & \tabularnewline
Dimanche & \multicolumn{6}{@{}p{11.0cm}}{\footnotesize Bilan et test de la semaine (30 min)} & \case & \tabularnewline
\hline
\multicolumn{9}{@{}p{16.2cm}@{}}{\footnotesize\textbf{Total de la semaine} \quad Jours étudiés : \trou[1.2cm] / 7 \quad Minutes : \trou[1.6cm] \quad Mots nouveaux : \trou[1.2cm] \quad Minutes de parole : \trou[1.4cm]} \tabularnewline
\end{tabularx}
\end{tableau}

\begin{tableau}[Les sept activités, en clair]
\footnotesize
\begin{tabularx}{\linewidth}{@{}p{2.4cm}p{1.4cm}X@{}}
\entete{Activité} & \entete{Durée} & \entete{Ce qu'elle contient} \tabularnewline
\hline
\textbf{Anki} & 15 min & 15 cartes nouvelles + révision du jour. Toujours avec une phrase entière. \tabularnewline
\textbf{Grammaire} & 15 min & Une unité de ce livre : règle, exemples lus à voix haute, un exercice corrigé. \tabularnewline
\textbf{Écoute} & 20 min & Un même audio, protocole des 5 passes (chapitre 16). Le même texte trois jours de suite. \tabularnewline
\textbf{Lecture} & 15 min & Un texte à 98\,\% de mots connus. Plaisir d'abord, dictionnaire ensuite. \tabularnewline
\textbf{Écrit} & 10 min & Journal : 80 à 100 mots sur une des 60 invites du chapitre 17. \tabularnewline
\textbf{Oral} & 20 min & Conversation réelle, ou monologue enregistré, ou un des 7 drills du chapitre 15. \tabularnewline
\textbf{Bonus} & --- & Chanson, série, film, message vocal à un ami, téléphone en anglais. \tabularnewline
\end{tabularx}
\end{tableau}

\section{Le registre annuel}
Une ligne par semaine. Ce tableau tient sur une seule page : affichez-le.

\begin{tableau}[52 semaines --- remplissez le dimanche soir]
\footnotesize
\begin{longtable}{@{}p{1.4cm}p{1.6cm}p{2.2cm}p{2.2cm}p{2.4cm}p{2.2cm}p{2.0cm}@{}}
\entete{Sem.} & \entete{Dates} & \entete{Jours /7} & \entete{Minutes totales} & \entete{Minutes de parole} & \entete{Mots actifs} & \entete{Humeur} \tabularnewline
\endfirsthead
\entete{Sem.} & \entete{Dates} & \entete{Jours} & \entete{Minutes} & \entete{Parole} & \entete{Mots} & \entete{Humeur} \tabularnewline
\endhead
\hline
1 & \trou[1.4cm] & \trou[2.0cm] & \trou[2.0cm] & \trou[2.2cm] & \trou[2.0cm] & \trou[1.8cm] \tabularnewline
2 & \trou[1.4cm] & \trou[2.0cm] & \trou[2.0cm] & \trou[2.2cm] & \trou[2.0cm] & \trou[1.8cm] \tabularnewline
3 & \trou[1.4cm] & \trou[2.0cm] & \trou[2.0cm] & \trou[2.2cm] & \trou[2.0cm] & \trou[1.8cm] \tabularnewline
4 & \trou[1.4cm] & \trou[2.0cm] & \trou[2.0cm] & \trou[2.2cm] & \trou[2.0cm] & \trou[1.8cm] \tabularnewline
5 & \trou[1.4cm] & \trou[2.0cm] & \trou[2.0cm] & \trou[2.2cm] & \trou[2.0cm] & \trou[1.8cm] \tabularnewline
6 & \trou[1.4cm] & \trou[2.0cm] & \trou[2.0cm] & \trou[2.2cm] & \trou[2.0cm] & \trou[1.8cm] \tabularnewline
7 & \trou[1.4cm] & \trou[2.0cm] & \trou[2.0cm] & \trou[2.2cm] & \trou[2.0cm] & \trou[1.8cm] \tabularnewline
8 & \trou[1.4cm] & \trou[2.0cm] & \trou[2.0cm] & \trou[2.2cm] & \trou[2.0cm] & \trou[1.8cm] \tabularnewline
9 & \trou[1.4cm] & \trou[2.0cm] & \trou[2.0cm] & \trou[2.2cm] & \trou[2.0cm] & \trou[1.8cm] \tabularnewline
10 & \trou[1.4cm] & \trou[2.0cm] & \trou[2.0cm] & \trou[2.2cm] & \trou[2.0cm] & \trou[1.8cm] \tabularnewline
11 & \trou[1.4cm] & \trou[2.0cm] & \trou[2.0cm] & \trou[2.2cm] & \trou[2.0cm] & \trou[1.8cm] \tabularnewline
12 & \trou[1.4cm] & \trou[2.0cm] & \trou[2.0cm] & \trou[2.2cm] & \trou[2.0cm] & \trou[1.8cm] \tabularnewline
13 & \trou[1.4cm] & \trou[2.0cm] & \trou[2.0cm] & \trou[2.2cm] & \trou[2.0cm] & \trou[1.8cm] \tabularnewline
14 & \trou[1.4cm] & \trou[2.0cm] & \trou[2.0cm] & \trou[2.2cm] & \trou[2.0cm] & \trou[1.8cm] \tabularnewline
15 & \trou[1.4cm] & \trou[2.0cm] & \trou[2.0cm] & \trou[2.2cm] & \trou[2.0cm] & \trou[1.8cm] \tabularnewline
16 & \trou[1.4cm] & \trou[2.0cm] & \trou[2.0cm] & \trou[2.2cm] & \trou[2.0cm] & \trou[1.8cm] \tabularnewline
17 & \trou[1.4cm] & \trou[2.0cm] & \trou[2.0cm] & \trou[2.2cm] & \trou[2.0cm] & \trou[1.8cm] \tabularnewline
18 & \trou[1.4cm] & \trou[2.0cm] & \trou[2.0cm] & \trou[2.2cm] & \trou[2.0cm] & \trou[1.8cm] \tabularnewline
19 & \trou[1.4cm] & \trou[2.0cm] & \trou[2.0cm] & \trou[2.2cm] & \trou[2.0cm] & \trou[1.8cm] \tabularnewline
20 & \trou[1.4cm] & \trou[2.0cm] & \trou[2.0cm] & \trou[2.2cm] & \trou[2.0cm] & \trou[1.8cm] \tabularnewline
21 & \trou[1.4cm] & \trou[2.0cm] & \trou[2.0cm] & \trou[2.2cm] & \trou[2.0cm] & \trou[1.8cm] \tabularnewline
22 & \trou[1.4cm] & \trou[2.0cm] & \trou[2.0cm] & \trou[2.2cm] & \trou[2.0cm] & \trou[1.8cm] \tabularnewline
23 & \trou[1.4cm] & \trou[2.0cm] & \trou[2.0cm] & \trou[2.2cm] & \trou[2.0cm] & \trou[1.8cm] \tabularnewline
24 & \trou[1.4cm] & \trou[2.0cm] & \trou[2.0cm] & \trou[2.2cm] & \trou[2.0cm] & \trou[1.8cm] \tabularnewline
25 & \trou[1.4cm] & \trou[2.0cm] & \trou[2.0cm] & \trou[2.2cm] & \trou[2.0cm] & \trou[1.8cm] \tabularnewline
26 & \trou[1.4cm] & \trou[2.0cm] & \trou[2.0cm] & \trou[2.2cm] & \trou[2.0cm] & \trou[1.8cm] \tabularnewline
27 & \trou[1.4cm] & \trou[2.0cm] & \trou[2.0cm] & \trou[2.2cm] & \trou[2.0cm] & \trou[1.8cm] \tabularnewline
28 & \trou[1.4cm] & \trou[2.0cm] & \trou[2.0cm] & \trou[2.2cm] & \trou[2.0cm] & \trou[1.8cm] \tabularnewline
29 & \trou[1.4cm] & \trou[2.0cm] & \trou[2.0cm] & \trou[2.2cm] & \trou[2.0cm] & \trou[1.8cm] \tabularnewline
30 & \trou[1.4cm] & \trou[2.0cm] & \trou[2.0cm] & \trou[2.2cm] & \trou[2.0cm] & \trou[1.8cm] \tabularnewline
31 & \trou[1.4cm] & \trou[2.0cm] & \trou[2.0cm] & \trou[2.2cm] & \trou[2.0cm] & \trou[1.8cm] \tabularnewline
32 & \trou[1.4cm] & \trou[2.0cm] & \trou[2.0cm] & \trou[2.2cm] & \trou[2.0cm] & \trou[1.8cm] \tabularnewline
33 & \trou[1.4cm] & \trou[2.0cm] & \trou[2.0cm] & \trou[2.2cm] & \trou[2.0cm] & \trou[1.8cm] \tabularnewline
34 & \trou[1.4cm] & \trou[2.0cm] & \trou[2.0cm] & \trou[2.2cm] & \trou[2.0cm] & \trou[1.8cm] \tabularnewline
35 & \trou[1.4cm] & \trou[2.0cm] & \trou[2.0cm] & \trou[2.2cm] & \trou[2.0cm] & \trou[1.8cm] \tabularnewline
36 & \trou[1.4cm] & \trou[2.0cm] & \trou[2.0cm] & \trou[2.2cm] & \trou[2.0cm] & \trou[1.8cm] \tabularnewline
37 & \trou[1.4cm] & \trou[2.0cm] & \trou[2.0cm] & \trou[2.2cm] & \trou[2.0cm] & \trou[1.8cm] \tabularnewline
38 & \trou[1.4cm] & \trou[2.0cm] & \trou[2.0cm] & \trou[2.2cm] & \trou[2.0cm] & \trou[1.8cm] \tabularnewline
39 & \trou[1.4cm] & \trou[2.0cm] & \trou[2.0cm] & \trou[2.2cm] & \trou[2.0cm] & \trou[1.8cm] \tabularnewline
40 & \trou[1.4cm] & \trou[2.0cm] & \trou[2.0cm] & \trou[2.2cm] & \trou[2.0cm] & \trou[1.8cm] \tabularnewline
41 & \trou[1.4cm] & \trou[2.0cm] & \trou[2.0cm] & \trou[2.2cm] & \trou[2.0cm] & \trou[1.8cm] \tabularnewline
42 & \trou[1.4cm] & \trou[2.0cm] & \trou[2.0cm] & \trou[2.2cm] & \trou[2.0cm] & \trou[1.8cm] \tabularnewline
43 & \trou[1.4cm] & \trou[2.0cm] & \trou[2.0cm] & \trou[2.2cm] & \trou[2.0cm] & \trou[1.8cm] \tabularnewline
44 & \trou[1.4cm] & \trou[2.0cm] & \trou[2.0cm] & \trou[2.2cm] & \trou[2.0cm] & \trou[1.8cm] \tabularnewline
45 & \trou[1.4cm] & \trou[2.0cm] & \trou[2.0cm] & \trou[2.2cm] & \trou[2.0cm] & \trou[1.8cm] \tabularnewline
46 & \trou[1.4cm] & \trou[2.0cm] & \trou[2.0cm] & \trou[2.2cm] & \trou[2.0cm] & \trou[1.8cm] \tabularnewline
47 & \trou[1.4cm] & \trou[2.0cm] & \trou[2.0cm] & \trou[2.2cm] & \trou[2.0cm] & \trou[1.8cm] \tabularnewline
48 & \trou[1.4cm] & \trou[2.0cm] & \trou[2.0cm] & \trou[2.2cm] & \trou[2.0cm] & \trou[1.8cm] \tabularnewline
49 & \trou[1.4cm] & \trou[2.0cm] & \trou[2.0cm] & \trou[2.2cm] & \trou[2.0cm] & \trou[1.8cm] \tabularnewline
50 & \trou[1.4cm] & \trou[2.0cm] & \trou[2.0cm] & \trou[2.2cm] & \trou[2.0cm] & \trou[1.8cm] \tabularnewline
51 & \trou[1.4cm] & \trou[2.0cm] & \trou[2.0cm] & \trou[2.2cm] & \trou[2.0cm] & \trou[1.8cm] \tabularnewline
52 & \trou[1.4cm] & \trou[2.0cm] & \trou[2.0cm] & \trou[2.2cm] & \trou[2.0cm] & \trou[1.8cm] \tabularnewline
\hline
\multicolumn{2}{@{}l}{\textbf{Total de l'année}} & \trou[1.4cm] & \trou[1.6cm] & \trou[1.8cm] & \trou[1.6cm] & \tabularnewline
\end{longtable}
\end{tableau}

\section{Le registre des conversations}
À remplir après chaque échange oral --- réel ou enregistré. C'est le tableau le plus
important de cette annexe : c'est celui qui mesure \textbf{ce que vous voulez vraiment}.

\begin{tableau}[Registre de conversation --- à photocopier]
\footnotesize
\begin{longtable}{@{}p{1.4cm}p{2.0cm}p{1.2cm}p{2.5cm}p{2.5cm}p{2.3cm}p{1.7cm}@{}}
\entete{Date} & \entete{Partenaire} & \entete{Durée} & \entete{Sujet} & \entete{Ce qui m'a bloqué} & \entete{3 phrases à retenir} & \entete{Note /5} \tabularnewline
\endfirsthead
\entete{Date} & \entete{Qui} & \entete{min} & \entete{Sujet} & \entete{Blocage} & \entete{À retenir} & \entete{Note} \tabularnewline
\endhead
\hline
 & & & & & & \tabularnewline
 & & & & & & \tabularnewline
 & & & & & & \tabularnewline
 & & & & & & \tabularnewline
 & & & & & & \tabularnewline
 & & & & & & \tabularnewline
 & & & & & & \tabularnewline
 & & & & & & \tabularnewline
 & & & & & & \tabularnewline
 & & & & & & \tabularnewline
 & & & & & & \tabularnewline
 & & & & & & \tabularnewline
 & & & & & & \tabularnewline
 & & & & & & \tabularnewline
 & & & & & & \tabularnewline
 & & & & & & \tabularnewline
 & & & & & & \tabularnewline
 & & & & & & \tabularnewline
 & & & & & & \tabularnewline
 & & & & & & \tabularnewline
\end{longtable}
\end{tableau}

\begin{astuce}[La note sur 5]
Utilisez la grille du chapitre 19, réduite à une seule note globale :
\textbf{1} = j'ai subi la conversation ; \textbf{2} = j'ai répondu, avec effort ;
\textbf{3} = j'ai tenu, en simplifiant ; \textbf{4} = j'ai été à l'aise ;
\textbf{5} = j'ai oublié que je parlais anglais. La moyenne de vos dix dernières notes est
votre vrai niveau --- bien plus fiable que n'importe quel test.
\end{astuce}

\section{Le registre des enregistrements}
Enregistrez-vous au moins une fois par semaine. Ce tableau garde la trace de vos progrès
auditifs --- les seuls que vous puissiez réécouter six mois plus tard avec stupéfaction.

\begin{tableau}[Mes enregistrements]
\footnotesize
\begin{tabularx}{\linewidth}{@{}p{1.6cm}p{2.2cm}p{3.0cm}p{1.6cm}p{2.2cm}X@{}}
\entete{\#} & \entete{Date} & \entete{Sujet / exercice} & \entete{Durée} & \entete{Pauses longues} & \entete{Ce que j'ai remarqué} \tabularnewline
\hline
1 & & Jour 1 --- présentation de soi & & & \tabularnewline
2 & & Semaine 4 & & & \tabularnewline
3 & & Semaine 8 & & & \tabularnewline
4 & & Semaine 12 & & & \tabularnewline
5 & & Semaine 16 & & & \tabularnewline
6 & & Semaine 20 & & & \tabularnewline
7 & & Semaine 24 & & & \tabularnewline
8 & & Semaine 28 & & & \tabularnewline
9 & & Semaine 32 & & & \tabularnewline
10 & & Semaine 36 & & & \tabularnewline
11 & & Semaine 40 & & & \tabularnewline
12 & & Semaine 44 & & & \tabularnewline
13 & & Semaine 48 & & & \tabularnewline
14 & & Jour 365 --- présentation de soi (même sujet qu'au jour 1) & & & \tabularnewline
\end{tabularx}
\end{tableau}

\begin{regle}[Le test du jour 1 et du jour 365]
Le jour 1, enregistrez-vous en train de vous présenter pendant trois minutes --- même sujet,
mêmes questions. Le jour 365, refaites exactement le même enregistrement. Écoutez-les à la
suite. Cette comparaison vaut tous les diplômes du monde : c'est la seule preuve irréfutable
de ce que 900 heures de travail ont produit.
\end{regle}

\section{Le suivi des 100 erreurs}
Une case par erreur du chapitre 18. Cochez-la quand vous l'avez \textbf{éliminée}, c'est-à-dire
quand vous ne la commettez plus pendant trois semaines consécutives.

\begin{tableau}[Tableau de chasse des erreurs 1 à 100]
\footnotesize
\begin{tabular}{@{}p{0.72cm}p{0.72cm}p{0.72cm}p{0.72cm}p{0.72cm}p{0.72cm}p{0.72cm}p{0.72cm}p{0.72cm}p{0.72cm}p{0.72cm}p{0.72cm}p{0.72cm}p{0.72cm}@{}}
\case & \case & \case & \case & \case & \case & \case & \case & \case & \case & \case & \case & \case & \case \\
1 & 2 & 3 & 4 & 5 & 6 & 7 & 8 & 9 & 10 & 11 & 12 & 13 & 14 \\[2pt]
\case & \case & \case & \case & \case & \case & \case & \case & \case & \case & \case & \case & \case & \case \\
15 & 16 & 17 & 18 & 19 & 20 & 21 & 22 & 23 & 24 & 25 & 26 & 27 & 28 \\[2pt]
\case & \case & \case & \case & \case & \case & \case & \case & \case & \case & \case & \case & \case & \case \\
29 & 30 & 31 & 32 & 33 & 34 & 35 & 36 & 37 & 38 & 39 & 40 & 41 & 42 \\[2pt]
\case & \case & \case & \case & \case & \case & \case & \case & \case & \case & \case & \case & \case & \case \\
43 & 44 & 45 & 46 & 47 & 48 & 49 & 50 & 51 & 52 & 53 & 54 & 55 & 56 \\[2pt]
\case & \case & \case & \case & \case & \case & \case & \case & \case & \case & \case & \case & \case & \case \\
57 & 58 & 59 & 60 & 61 & 62 & 63 & 64 & 65 & 66 & 67 & 68 & 69 & 70 \\[2pt]
\case & \case & \case & \case & \case & \case & \case & \case & \case & \case & \case & \case & \case & \case \\
71 & 72 & 73 & 74 & 75 & 76 & 77 & 78 & 79 & 80 & 81 & 82 & 83 & 84 \\[2pt]
\case & \case & \case & \case & \case & \case & \case & \case & \case & \case & \case & \case & \case & \case \\
85 & 86 & 87 & 88 & 89 & 90 & 91 & 92 & 93 & 94 & 95 & 96 & 97 & 98 \\[2pt]
\case & \case & & & & & & & & & & & & \\
99 & 100 & & & & & & & & & & & & \\
\end{tabular}
\end{tableau}

\section{Le bilan mensuel}
À photocopier douze fois. Dix minutes, le dernier dimanche de chaque mois.

\begin{tableau}[Mois \rule{1cm}{0.4pt} --- du \rule{2cm}{0.4pt} au \rule{2cm}{0.4pt}]
\footnotesize
\begin{tabularx}{\linewidth}{@{}p{7.6cm}X@{}}
\entete{Question} & \entete{Réponse} \tabularnewline
\hline
Nombre de jours étudiés sur 30 & \rule{0pt}{0.9cm}\trou[5.0cm] \tabularnewline
Nombre total de minutes & \rule{0pt}{0.9cm}\trou[5.0cm] \tabularnewline
Nombre de minutes de parole réelle & \rule{0pt}{0.9cm}\trou[5.0cm] \tabularnewline
Nombre de conversations tenues & \rule{0pt}{0.9cm}\trou[5.0cm] \tabularnewline
Nombre de mots actifs dans Anki & \rule{0pt}{0.9cm}\trou[5.0cm] \tabularnewline
Chapitres de ce livre terminés & \rule{0pt}{0.9cm}\trou[5.0cm] \tabularnewline
Test mensuel --- score de grammaire & \rule{0pt}{0.9cm}\trou[2.0cm] / 20 \tabularnewline
Test mensuel --- score de vocabulaire & \rule{0pt}{0.9cm}\trou[2.0cm] / 30 \tabularnewline
Grille orale --- total & \rule{0pt}{0.9cm}\trou[2.0cm] / 25 \tabularnewline
Erreurs éliminées ce mois-ci (numéros) & \rule{0pt}{0.9cm}\trou[5.0cm] \tabularnewline
Ce qui a le mieux marché & \rule{0pt}{1.4cm}\trou[6.0cm] \tabularnewline
Ce qui m'a fait perdre du temps & \rule{0pt}{1.4cm}\trou[6.0cm] \tabularnewline
Ce que je change le mois prochain & \rule{0pt}{1.4cm}\trou[6.0cm] \tabularnewline
Mon moment de fierté du mois & \rule{0pt}{1.4cm}\trou[6.0cm] \tabularnewline
Un mot d'encouragement pour moi-même & \rule{0pt}{1.4cm}\trou[6.0cm] \tabularnewline
\end{tabularx}
\end{tableau}

\section{La lettre du jour 365}
Le dernier jour, écrivez-vous une lettre --- \textbf{en anglais}, évidemment. Elle répond
à quatre questions : qu'est-ce que je sais faire aujourd'hui que je ne savais pas faire il y
a un an ? Quelle est la conversation dont je suis le plus fier ? Qu'est-ce qui m'a fait
vouloir abandonner, et comment ai-je continué ? Que vais-je faire de cette compétence ?
Puis cachetez-la, et ouvrez-la dans cinq ans.

\begin{tableau}[Mes repères de fin d'année]
\begin{tabularx}{\linewidth}{@{}p{8.2cm}X@{}}
\entete{Indicateur} & \entete{Jour 1 \qquad\qquad Jour 365} \tabularnewline
\hline
Je peux me présenter et parler de ma vie pendant\ldots{} & \trou[2.2cm] min \qquad \trou[2.2cm] min \tabularnewline
Je comprends un anglophone qui parle normalement\ldots{} & \trou[2.2cm] \% \qquad \trou[2.2cm] \% \tabularnewline
Je connais activement\ldots{} & \trou[2.2cm] mots \quad \trou[2.2cm] mots \tabularnewline
Je lis un article de presse\ldots{} & \trou[2.2cm] \% \qquad \trou[2.2cm] \% \tabularnewline
J'écris un e-mail correct en\ldots{} & \trou[2.2cm] min \qquad \trou[2.2cm] min \tabularnewline
Je traduis encore dans ma tête\ldots{} & \trou[2.2cm] fois/min \quad \trou[2.2cm] fois/min \tabularnewline
Test de positionnement (chapitre 19)\ldots{} & \trou[2.2cm] /40 \qquad \trou[2.2cm] /40 \tabularnewline
Grille d'auto-évaluation orale\ldots{} & \trou[2.2cm] /25 \qquad \trou[2.2cm] /25 \tabularnewline
\end{tabularx}
\end{tableau}

\begin{mission}[Commencez maintenant, pas demain]
Imprimez cette annexe aujourd'hui. Remplissez le contrat personnel ce soir --- même
imparfaitement, même en dix minutes. Fixez l'heure de votre première séance, et mettez une
alarme sur votre téléphone pour demain matin. Le jour 1 n'est pas une date dans un livre :
c'est le moment où vous prenez le stylo.
\end{mission}
